\SourceLeaf{151}hang, und erkenne darin nichts weiter als eine baare Petitio principii. Sie werden mich nicht missverstehen. Ich lege wenig Werth auf eine streng logische Einkleidung, die, in Schriften für Männer und Kenner, oft nur Pedanterie sein würde; aber der streng logische Zusammenhang muss sich, wo es nur gefordert wird, überall nachweisen lassen, diess ist eine unerlässliche Bedingung eines guten Vortrags. Wenn ich einen Mathematiker sehr hochschätze, und hierin etwas vermisse, so denke ich immer zuerst, dass es \textls{bloss} im Vortrag liegt, und gehe schwer daran, anzunehmen, dass der Gedanke selbst leer ist. Aber bei Herrn \textsc{Ivorys} drei Aufsätzen bin ich leider dazu genöthigt.

Die Abhandlung über die Kleinsten Quadrate ist denn doch wirklich unter aller Kritik. Welche Verworrenheit, Unklarheit und völliger Mangel logischer Bündigkeit! Aber darüber könnte man noch wegsehen, wenn wirklich ein reeller Gedanke zum Grunde läge. Allein das ist durchaus nicht der Fall. Wie ist es möglich, dass das Geschwätz p. 163 und 164 uns als ein \textls{proof} that this principle leads \textls{necessarily} to the method of the least squares venditirt wird!

In der That weit entfernt, dass daraus eine \textls{Nothwendigkeit} der M. der kl. Q. folge, kann man dieses Ganze\HistoricalNote{*)}{*)}{Versteht sich mit Ausnahme der Folge, dass $\varepsilon\varepsilon+\varepsilon'\varepsilon'+\varepsilon''\varepsilon''+\cdots$ ein Minimum wird; aber diess soll ja eben nicht \textls{Axiom} sein, sondern bewiesen werden.} auf \textls{jede beliebige} andere Behandlung der vorgegebenen Gleichungen anwenden; z. B. wenn, anstatt sie mit $a$, $a'$, $a''$ zu multipliciren, man sie mit irgend einer Potenz dieser Coefficienten multiplicirte, oder auch sie dividirte. Die Abhandlung scheint in der That halb im Schlafe geschrieben zu sein, und Sie erlassen mir wohl noch mehrere einzelne Stellen anzustechen, die Sie sogleich selbst bemerken werden, sobald Sie nur nicht von einer vorgefassten Meinung ausgehen, dass irgend etwas daran ist.

Wenig besser scheint mir der andere Aufsatz über die Pendellängen; die Vorwürfe, die er der Anwendung der M. der kl. Q. macht, zeigen nur, dass ihr Geist dem Hrn. I. ganz fremd ist. Ich finde, dass Hr. \textsc{Sabine} diese Methode im Wesentlichen ganz richtig angewandt hat (— die Form der Anwendung ist allerdings nicht die zweckmässigste, und ich habe, unter uns gesagt, eigentlich \textsc{Sabine}, \textsc{Paucker} und \textsc{Muncke} in der neuen Ausgabe von \textsc{Gehler}, bei meiner
